% ============================================================
%  吴万星 — 个人简历 (中文版)
% ============================================================
\documentclass[10pt, a4paper]{article}

% ── Encoding & CJK Support ──────────────────────────────────
\usepackage[UTF8, heading=false]{ctex}
\usepackage[T1]{fontenc}
\usepackage{palatino}
\usepackage{mathpazo}
\usepackage{microtype}

% ── Page Layout ──────────────────────────────────────────────
\usepackage[top=1.2cm, bottom=1.2cm, left=1.5cm, right=1.5cm]{geometry}
\setlength{\parindent}{0pt}
\setlength{\parskip}{0pt}

% ── Packages ─────────────────────────────────────────────────
\usepackage{xcolor}
\usepackage{hyperref}
\usepackage{enumitem}
\usepackage{titlesec}
\usepackage{tabularx}
\usepackage{array}
\usepackage{fancyhdr}
\usepackage{tikz}
\usepackage{etoolbox}
\usepackage{graphicx}

% ── Colors ───────────────────────────────────────────────────
\definecolor{headcol}{RGB}{30, 30, 30}
\definecolor{rulecol}{RGB}{180, 50, 30}
\definecolor{secline}{RGB}{200, 200, 200}
\definecolor{linkcol}{RGB}{40, 80, 160}
\definecolor{dategray}{RGB}{100, 100, 100}

% Tag colors
\definecolor{tagroute}{RGB}{130, 50, 160}
\definecolor{taggenomics}{RGB}{40, 140, 90}
\definecolor{tagrl}{RGB}{180, 100, 20}

\hypersetup{
  colorlinks = true,
  urlcolor   = linkcol,
  linkcolor  = linkcol,
  pdfauthor  = {吴万星},
  pdftitle   = {吴万星 — 个人简历},
}

% ── Section Headers ──────────────────────────────────────────
\titleformat{\section}
  {\large\bfseries\color{headcol}}
  {}{0em}{}
  [\vspace{2pt}{\color{secline}\rule{\linewidth}{0.4pt}}\vspace{-2pt}]
\titlespacing*{\section}{0pt}{8pt}{3pt}

% ── Badge Tags ───────────────────────────────────────────────
\newcommand{\badge}[2]{%
  \tikz[baseline=(x.base)]{
    \node[inner xsep=3.5pt, inner ysep=1.6pt, rounded corners=2.5pt,
          draw=#1, line width=0.45pt,
          text=#1, font=\fontsize{6}{7}\selectfont\sffamily]
      (x) {#2};
  }%
}

\newcommand{\troute}  {\badge{tagroute}{routing}}
\newcommand{\tgenomics}{\badge{taggenomics}{genomics}}
\newcommand{\trl}     {\badge{tagrl}{RL}}

% ── Bullet Lists ─────────────────────────────────────────────
\setlist[itemize,1]{
  leftmargin = 1.3em,
  label      = {\small\textbullet},
  itemsep    = 1pt,
  topsep     = 2pt,
  parsep     = 0pt,
}

% ── Helpers ──────────────────────────────────────────────────
\newcommand{\unilogo}[1]{%
  \raisebox{-0.5pt}{\includegraphics[height=8pt]{assets/logos/cv/#1}}%
}
\newcommand{\me}[1]{\textbf{#1}}
\newcommand{\eqstar}{$^{*}$}
\newcommand{\corr}{$^{\dagger}$}

\def\pubskip{3pt}

% Single publication block
\newcommand{\pub}[4]{%
  \vspace{\pubskip}%
  \noindent\textbf{#1}\ifstrempty{#2}{}{\enskip #2}%
  \begin{itemize}[topsep=0pt]
    \item #3
    \item \textit{#4}
  \end{itemize}%
}

% Education / Experience entry
\newcommand{\entry}[2]{%
  \noindent
  \begin{tabularx}{\linewidth}{@{}X r@{}}
    \textbf{#1} & {\small\color{dategray}#2}
  \end{tabularx}%
  \par\vspace{0pt}%
}

% ── Header / Footer ──────────────────────────────────────────
\pagestyle{fancy}
\fancyhf{}
\renewcommand{\headrulewidth}{0pt}
\fancyfoot[R]{\footnotesize\color{dategray}\thepage}

% ============================================================
\begin{document}
\thispagestyle{empty}

% ── NAME BLOCK ───────────────────────────────────────────────
\noindent
\begin{tabularx}{\linewidth}{@{}X r@{}}
  \raisebox{-4pt}{%
    {\fontsize{28}{32}\selectfont\bfseries 吴万星}%
    \quad{\large\color{dategray}2027届}%
  }
  &
  \begin{tabular}[b]{@{}r@{}}
    \small\href{mailto:wuwanxing3574@gmail.com}{wuwanxing3574@gmail.com}\\[2pt]
    \small\href{https://github.com/wwx777}{GitHub}
    \quad \href{https://wwx777.github.io}{个人主页}
  \end{tabular}
\end{tabularx}

\vspace{6pt}
{\color{rulecol}\rule{\linewidth}{1.2pt}}
\vspace{4pt}

% ── ABOUT ────────────────────────────────────────────────────
\section{个人简介}

\textbf{巴黎综合理工学院（Institut Polytechnique de Paris）}计算机科学硕士在读（2025--2027），
本科毕业于\textbf{南方科技大学}数据科学与大数据技术专业。
现为\textbf{SUSTech-NLP}（陈冠华教授）科研助理，
研究方向为\textbf{大语言模型协作路由系统}，
聚焦于 LLM--SLM 路由的忠实性、泛化性与鲁棒性。

\vspace{3pt}
主要工作包括：路由公平性评测基准 RouterXBench 的设计，
以及基于 Dirichlet 的动态隐层探针路由方法。
一作论文 1 篇（在审）。

% ── EDUCATION ────────────────────────────────────────────────
\section{教育经历}

\entry{\unilogo{ipp.png}\enskip 巴黎综合理工学院（Institut Polytechnique de Paris）}{法国巴黎}
\textit{计算机科学硕士} \hfill \textit{2025.09 -- 2027.07}

\vspace{4pt}
\entry{\unilogo{sustech.png}\enskip 南方科技大学}{深圳}
\textit{数据科学与大数据技术，工学学士} \hfill \textit{2021.09 -- 2025.07}
\begin{itemize}
  \item 智仁书院优秀毕业生（2025）
\end{itemize}

% ── EXPERIENCE ───────────────────────────────────────────────
\section{科研经历}

\entry{\unilogo{sustech.png}\enskip SUSTech-NLP 实验室}{深圳}
\textit{科研助理}，导师：陈冠华教授 \hfill \textit{2025.01 -- 至今}
\begin{itemize}
  \item \textbf{RouterXBench}：设计 LLM--SLM 路由分布偏移忠实性评测矩阵。
  \item \textbf{动态隐层探针}：开发基于 Dirichlet 分布的分布感知逐层聚合路由方法。
\end{itemize}

\section{实习经历}

\entry{华大基因（BGI）智能医学研究院（IIMR）}{深圳}
\textit{机器学习实习生（AI for Genomics）} \hfill \textit{2024.07 -- 2025.01}
\begin{itemize}
  \item 构建 cfDNA 特征表示（统计筛选 + 模型导向特征选择），集成学习用于结直肠癌早筛，
        实现 \textbf{94.88\% 灵敏度}。
\end{itemize}

\vspace{4pt}
\entry{\unilogo{pku.png}\enskip 北京大学深圳医院}{深圳}
\textit{数据分析实习生} \hfill \textit{2023.07 -- 2023.10}
\begin{itemize}
  \item 对多时间点临床记录进行生物统计分析，比较不同手术方式的术后结局。
\end{itemize}

% ── PUBLICATIONS ─────────────────────────────────────────────
\section{学术论文}

\pub{Towards Fair and Comprehensive Evaluation of Routers in Collaborative LLM Systems}{\troute}
    {\me{Wanxing Wu}\eqstar, He Zhu\eqstar, Yixia Li\eqstar, Yun Chen, Guanhua Chen\corr
     \quad ($^*$同等贡献)}
    {在审，2026 \quad[\href{https://arxiv.org/abs/2602.11877}{arXiv:2602.11877}]}

% ── PROJECTS ─────────────────────────────────────────────────
\section{项目经历}

\noindent\textbf{ReflexSQL：自反思 Text-to-SQL Agent} \hfill \textit{2025.12 -- 2026.01}
\begin{itemize}
  \item 在 Spider~2.0-Lite 上设计自反思 LLM agent，结合查询内反思修正循环与跨查询经验记忆缓冲区。
  \item Qwen3-32B 骨干模型可执行准确率从 16.29\%（基线）提升至 \textbf{42.53\%}；
        完成骨干模型对比与消融实验。
\end{itemize}

\vspace{4pt}
\noindent\textbf{保守离策略置信域策略优化（Conservative Off-Policy TRPO）} \hfill \textit{2025.11 -- 2026.01}
\begin{itemize}
  \item 针对离策略 TRPO 设计逐维掩码策略，利用在策略梯度对齐抑制冲突更新分量，
        确保更新方向与在策略目标一致；在 PyBullet/Gym 运动控制任务上验证。
\end{itemize}

% ── HONORS ───────────────────────────────────────────────────
\section{荣誉与奖项}

\begin{itemize}[itemsep=0pt, topsep=1pt]
  \item \textbf{优秀毕业生}，南方科技大学智仁书院（2025.07）
  \item 优秀学生奖学金，南方科技大学（2024.10）
  \item \textbf{美国大学生数学建模竞赛（MCM/ICM）}，Honorable Mention（2024.02）
\end{itemize}

% ── SKILLS ───────────────────────────────────────────────────
\section{技能}

\noindent\textbf{大模型与机器学习：} 推理框架（vLLM、SGLang），Agentic-RL 工作流，SFT/RL 训练，模型评测。\\[3pt]
\noindent\textbf{编程与工具：} Python、PyTorch、Docker、Linux、Git、Shell 脚本。\\[3pt]
\noindent\textbf{语言：} 英语（IELTS 6.5，CET-4，CET-6）；中文（母语）。

\end{document}
